\documentclass{article}
\usepackage[T1]{fontenc}
\usepackage{lmodern}
\usepackage{graphicx}
\usepackage{amsmath,amssymb}
\usepackage[a4paper,margin=2cm]{geometry}
\usepackage[absolute,overlay]{textpos}
\setlength{\TPHorizModule}{1mm}
\setlength{\TPVertModule}{1mm}
\usepackage[indonesian]{babel}
\usepackage{hyperref}
\setlength{\emergencystretch}{2em}
% unit: Halaman 12

\begin{document}

% === Halaman 12 (python/native) ===

12

Cipher Sederhana dengan operasi XOR

• Sama seperti Vigenere Cipher, tetapi dalam mode bit

• Setiap bit plainteks di-XOR-kan dengan setiap bit kunci.


Enkripsi: C =  P  K 


Dekripsi: P = C  K  


Contoh: 

plainteks 


01100101  

(karakter ‘e’) 



kunci  


00110101  

(karakter ‘5’) 




cipherteks  

01010000  

(karakter ‘P’) 



kunci  


00110101  

(karakter ‘5’) 




plainteks 


01100101  

(karakter ‘e’)

\end{document}
